Durante o estágio realizado na Veeries, entre outubro de 2025 e abril de 2026, identificou-se a necessidade de reestruturar a ingestão de dados de fontes externas críticas para o negócio (confidencial). 
O problema central era a ausência de um fluxo de dados padronizado, o que comprometia a rastreabilidade e a confiabilidade das informações consumidas pelas áreas de análise.
Para solucionar essa questão, atuei diretamente no desenvolvimento de um pipeline de dados automatizado, aplicando fundamentos de automação de processos do curso de Engenharia de Controle e Automação.

\section*{Objetivos e Metodologia}

O objetivo geral do projeto foi desenvolver um fluxo robusto de extração, processamento e disponibilização de dados.
Como objetivos específicos, destaco a construção do fluxo baseada na Arquitetura Medallion (camadas Bronze, Silver e Gold) e a implementação de rigorosas validações de qualidade estrutural.
\section*{Estrutura do Relatório}
Para relatar o desenvolvimento e os resultados alcançados, este documento está organizado da seguinte forma:
\begin{itemize}
    \item Capítulo 2: Descreve o contexto da empresa e detalha o problema enfrentado com os dados.
    \item Capítulo 3: Apresenta brevemente os conceitos da Arquitetura Medallion e as técnicas de validação utilizadas.
    \item Capítulo 4: Expõe a implementação técnica da solução, detalhando a codificação do pipeline (L0 a L2).
    \item Capítulo 5 e 6: Apresentam os resultados obtidos, a análise crítica dos impactos gerados e as considerações finais.
\end{itemize}


\begin{itemize}
\item 	definição do assunto, o problema a ser abordado no trabalho;
\item 	localização do assunto no espaço, no tempo e no contexto do Curso;
\item 	importância do assunto e justificativa da escolha;
\item 	objeto geral e objetivos específicos;
\item 	hipóteses levantadas ou argumentação principal;
\item 	descrição da metodologia empregada no trabalho (bibliografia e/ou estudo de campo e/ou de laboratório);
\item 	as alternativas de solução do Problema, decorrentes de estudo bibliográfico que deverá ficar evidenciado mediante referências;
\item 	apresentação do plano adotado para o desenvolvimento do assunto, isto é, a descrição do que será tratado em cada um dos capítulos subsequentes.
\end{itemize}

